\documentclass[11pt]{article}
\usepackage[margin=1in]{geometry}
\usepackage{hyperref}
\usepackage{microtype}
\hypersetup{colorlinks=true,linkcolor=blue,citecolor=blue,urlcolor=blue}
\title{A Research-Use-Only Healthspan Assessment Engine with Explicit Evidence Boundaries\\\large System design, deterministic deficit accumulation, and clinical validation gates}
\author{Stan Chen}
\date{23 November 2026\\[4pt]\small Unpublished working draft; not submitted to arXiv}
\begin{document}
\maketitle
\begin{center}\small This draft has no arXiv identifier. Claims are limited to the evidence described below.\end{center}
\begin{abstract}
This paper describes research-use-only software for organizing healthspan assessment data and presenting a deterministic deficit-accumulation index with a guarded biological-age interface. The engine accepts a 35-feature assessment, enforces a minimum viable vector, and computes a Rockwood-style frailty index. A synthetic browser companion demonstrates local review and follow-up workflows. The repository explicitly does not establish clinical approval, a validated survival model, a clinical reference panel, or validated uncertainty analysis; it makes no lifespan, diagnostic, mortality, or treatment-effect claim. No patient cohort or clinical outcome evaluation is reported. We describe system boundaries and a validation program requiring governance, representative data, external validation, uncertainty analysis, and workflow evaluation.
\end{abstract}
\section{Intended use}
Healthspan software can appear clinically authoritative before evidence supports that use. This project treats intended use and evidence boundaries as product requirements. The current software is a research prototype and evidence showcase, not a clinical decision system.\par
Deficit accumulation is an established research construct, but an implementation does not validate a feature set or resulting score for a population or clinical decision \cite{r1}.\par
\section{System design}
The Python engine validates a 35-feature assessment and computes a deterministic Rockwood-style index; a biological-age interface is guarded. A synthetic clinic companion demonstrates save, review, approval, check-in, and follow-up export. It has no patient accounts, clinic authentication, or hosted patient records.\par
Local manual and CSV flows are distinct from PDF extraction, which may send de-identified page content to an external AI service only after explicit consent. A browser preview or reproducible calculation cannot establish clinical validity, utility, or regulatory status.\par
\section{Evidence and validation gates}
The repository states \cite{r3} that it does not ship an approved survival model, a clinical reference panel, or validated uncertainty analysis, and that E-005 remains blocked. Software acceptance and scientific evidence are distinct.\par
Before clinical claims, fix intended population and measurement definitions; assess missingness, reliability, calibration, subgroup performance, uncertainty, external transportability, and decision consequences. Prospective work requires appropriate ethics review and data governance. Reporting should follow relevant prediction-model guidance such as TRIPOD+AI \cite{r2}.\par
\section{Limitations and conclusion}
This draft reports no cohort, outcome, calibration, survival, treatment-effect, or prospective workflow results. Deterministic behavior does not establish clinical validity.\par
The current contribution is a bounded prototype with explicit evidence gaps. Clinical utility requires governance, representative validation, uncertainty analysis, and workflow study.\par
\begin{thebibliography}{99}
\bibitem{r1} Rockwood, K., and Mitnitski, A. Frailty in relation to the accumulation of deficits. Journal of Gerontology: Medical Sciences, 62(7), 722-727, 2007. \url{https://doi.org/10.1093/gerona/62.7.722}
\bibitem{r2} Collins, G. S., et al. TRIPOD+AI statement. BMJ, 2024. \url{https://www.bmj.com/content/385/bmj-2023-078378}
\bibitem{r3} Chen, S. Clinical Healthspan and Deficit Accumulation Engine. \url{https://github.com/stancsz/healthspan-review}
\end{thebibliography}
\end{document}
